% ======================================================================
\section{Checks and limits}\label{sec:checks}
% ======================================================================

The calculation involves a large number of terms, with non-commuting charges and Wilson lines in a definite order. In this section we list the checks to which the result has been subjected, with what each of them tests and what it does not, and compare the result with the literature. What has not been checked is stated explicitly, and collected in Section~\ref{sec:openpoints}.

\subsection{Overview of the checks}\label{sec:checks_table}

Table~\ref{tab:checks} lists the checks.
\begingroup
\footnotesize
\renewcommand{\arraystretch}{1.25}
\begin{longtable}{@{}>{\raggedright\arraybackslash}p{0.27\textwidth}>{\raggedright\arraybackslash}p{0.14\textwidth}>{\raggedright\arraybackslash}p{0.25\textwidth}>{\raggedright\arraybackslash}p{0.25\textwidth}@{}}
\caption{The checks of the calculation. ``A'': analytic; ``N'': numerical, with the accuracy reached.}\label{tab:checks}\\
\toprule
check & kind, accuracy & what it tests & what it does not test\\
\midrule
\endfirsthead
\toprule
check & kind, accuracy & what it tests & what it does not test\\
\midrule
\endhead
\bottomrule
\endfoot
Perturbation series with operator-valued matrix elements, Eqs.~\eqref{eq:pt2}, \eqref{eq:pt3} (App.~\ref{app:pt_num}) & N, scaling $g^{N+1}$ & order of the factors, normalisation terms & continuum matrix elements\\
Components of the dressed one-gluon state (Sec.~\ref{sec:onegin}) & N & $\mf B_3$, $\mf C_1$ from the matrix elements, signs, double pole & coordinate-space forms\\
Unitarity relations (App.~\ref{app:unitarity}) & A, symbolic & all relations to $O(g^3)$ with non-commuting coefficients & the explicit coefficients\\
Amplitudes~\eqref{eq:A3parts}, \eqref{eq:A2parts} (App.~\ref{app:rows_check}) & A, symbolic; $4\times10^{-16}$ & signs, factors, order of the coefficients, Wilson lines & factors $(2\pi)^n$, transverse arguments\\
Principal-value Fourier transforms~\eqref{eq:FT_i3}--\eqref{eq:FT_i5} (App.~\ref{app:fourier}) & N, $10^{-15}$ & the transforms as distributions & ---\\
Ordering $[\mf A^{(3)}]_{\rho\rho\rho}=\mf A\mf N_2$ against an exact construction of $\Omega$ with matrices (Sec.~\ref{sec:check_wf}) & N, $3\times10^{-16}$ & operator order in third-order LCPT & the transverse and longitudinal integrals\\
Wave function against Lublinsky--Mulian (Sec.~\ref{sec:check_wf}) & A & all components to $O(g^3)$, incl.\ $\frac{11}{3}$, $\frac{67}{9}-\frac{\pi^2}{3}$, cut-off logarithms & operator order of the three-$\rho$ term; coordinate-space forms\\
$U=1$, constant $U$ (Sec.~\ref{sec:limits}) & A, manifest & every term is built from differences & relative coefficients of terms\\
Four-$\rho$ cancellation (Sec.~\ref{sec:fourrho}) & A; N with $SU(2)$, $SU(3)$ matrices & supports of $\mf B_2$, normalisation terms, ordering & two- and three-$\rho$ kernels\\
Symmetrised three-$\rho$ logarithmic terms vanish (Sec.~\ref{sec:threerho}) & A; N, $2\times10^{-15}$ & colour tensors and kernels of all cubic logarithmic terms, with contact terms & non-logarithmic cubic terms\\
BFKL list $=$ Eq.~\eqref{eq:BFKLrho} $=$ singlet form (Sec.~\ref{sec:check_BFKL}) & N, $10^{-14}$ & regrouping and colour reduction & the list itself\\
BFKL term against exact one-gluon evolution of a dipole (Sec.~\ref{sec:check_BFKL}) & A, pointwise in $\v z$ & the list itself, sign and normalisation, at finite $N_c$, for every colour-singlet projectile & non-singlet projectiles\\
JIMWLK list $=-H_{\rm JIMWLK}$ on LO (Sec.~\ref{sec:check_JIMWLK}) & N, $10^{-14}$ & all kernels, signs, normalisation, for an arbitrary ordered correlator & finite terms\\
Double logarithms: $\sum T_{\rm row}=0$ (App.~\ref{sec:fin_soft}) & N, $<10^{-7}$ vs.\ $O(1)$ & relative weights of all rows with recoil & terms without a tail\\
Pairs A--D: colour identities and tails (App.~\ref{sec:fin_soft}) & A; N, $2\times10^{-15}$; sum $O(|\v z|^{-3})$ & that soft remainders cancel for any $U(\v z)$ & absolute normalisation of a pair\\
Recoil identity of Row~I.III, Eq.~\eqref{eq:recoil_identity} & N, $3\times10^{-13}$ & exchange of $\Lambda\to0$ with the transverse integral & ---\\
Row~I.III after the $\v w$ integral, Eq.~\eqref{eq:Mhatclaim} & A (Eq.~\eqref{eq:sawid}); N at $|\v z|=2560$ & absence of a tail in all directions & a full transverse integration\\
Row~I.IV $+$ Row~III.I $=0$ (Sec.~\ref{sec:fin_together}) & A; N, $10^{-14}$ rel. & colour, kernel and phase of both rows & other rows\\
Pole balance (Table~\ref{tab:poles}) & A, exact rational arithmetic & all pole coefficients & finite parts\\
Identity~\eqref{eq:dg_SEid}, point-like projectile & N, $10^{-13}$, $\xi$ by $\xi$ & cancellation of the self-energy by the unresolved pair & resolved part of Row~I.V\\
$\MS$ counterterm: $e^+e^-\to gX$, $H\to gX$ coefficient functions (Sec.~\ref{sec:dg_msbar}) & A, exact & the Jacobian $\xi^{-d}$ of the fragmentation counterterm & the transverse-momentum dependence (angle-integrated checks)\\
Tail of Row~II.IV: weights $1:1:-2$ (Sec.~\ref{sec:dg_tail}) & A; N, $10^{-14}$ (colour) & cancellation of the soft-parent region between three rows & interchange of the $p^+$ integral with the limit in Row~III.III\\
Dilute limit, real emission $=$ QMRK two-gluon production (Sec.~\ref{sec:check_dilute}) & N, $2\times10^{-15}$ & amplitudes $\mathcal A_B$, $\mathcal A_{AA}$, $\mathcal A_{CA}$, contact parts, normalisation, $d_\perp$ dependence & the printed coordinate-space finite terms; virtual part\\
Dilute limit, virtual part, point-like quark (Sec.~\ref{sec:check_dilute}) & N, $2\times10^{-4}$ (logs), $2\times10^{-3}$ (constant) & Regge trajectories as coefficients of both logarithms; absence of $\frac{67}{9}$ & finite part against the one-loop vertex (open)\\
Cubic terms on a quark and a dipole (Sec.~\ref{sec:check_cubic}) & N, $5\times10^{-14}$ & the split~\eqref{eq:cubicavg} into $\sigma$ and contact parts; C-parity & kernels\\
Printed equations against verified colour tensors (Sec.~\ref{sec:lit_finite}) & N, $7\times10^{-14}$ & that the printed colour factors and averages are those of the calculation & kernels beyond the transcription\\
Integrability: power counting of all kernels (App.~\ref{sec:fin_integrable}) & A and N (exponent fits) & local integrability at all coincidences & every degenerate configuration; a full numerical integration\\
\end{longtable}
\endgroup

Unless stated otherwise, the algebraic identities quoted below were verified both analytically and numerically: the expressions were evaluated on a transverse lattice with random $SU(2)$ and $SU(3)$ Wilson lines in the adjoint representation and with either a completely general ordered charge correlator $\langle\rho^b(\v u)\rho^c(\v v)\rangle$ or the colour-singlet one, Eq.~\eqref{eq:singlet}. Since the Wilson lines are random, an identity that holds on the lattice to round-off error cannot hold by accident. The scripts used for the checks are available from the authors on request.

\subsection{Limits}\label{sec:limits}

\paragraph{No target, pure-gauge target.} For $U=1$ the operator $\bar\Omega^\dagger\Omega$ is the identity and every amplitude vanishes. In the result this is manifest: every term of every row, and every term of the BFKL, JIMWLK, DGLAP and running-coupling contributions, contains a factor that vanishes at $U=1$, such as $[U(\v x')-U(\v w')]^{ab'}\rho^{b'}(\v x')$, its analogue with three Wilson lines, or a combination of $f^{abc}$ and Wilson lines that reduces to $N_c-N_c-N_c+N_c=0$. A constant $U$ is a global colour rotation, which cannot produce gluons; all terms vanish for constant $U$ as well, by $f^{abc}U^{bd}U^{ce}=U^{ah}f^{hde}$ and $U^TU=1$.

\paragraph{Counting of charges.} The BFKL and JIMWLK terms are of order $g^4\rho^2$. No term of order $g^4\rho^4$ survives (Section~\ref{sec:fourrho}), and the three-$\rho$ terms that multiply a rapidity logarithm cancel in the symmetrised sense (Section~\ref{sec:threerho}). The finite terms are of order $g^4\rho^2$ and $g^4\rho^3$. After the average over the projectile the $g^4\rho^3$ terms do not drop out: their symmetrised part is proportional to the symmetrised three-point function, i.e.\ to the three-point function $\sigma$ of the projectile, Eq.~\eqref{eq:A2d}, and the commutators produced by the reordering of the charges leave contact terms of order $g^4\langle\rho\rho\rangle$, of the same order as the quadratic finite terms (Section~\ref{sec:fin_singlet}).

\paragraph{Reality.} Groups I and III are written with their complex conjugates and are real. Group II is Hermitian by itself: its kernels are symmetric under $(\v x,\v w)\leftrightarrow(\v x',\v w')$ and so is its colour operator, which we verified numerically to $10^{-12}$.

\paragraph{Odd terms.} The three-$\rho$ terms carry an explicit factor $i$ (in the cubic part of Row~I.II it is in the Fourier transform of an odd kernel), and all kernels depend on $\v k$ through functions with $f(-s)=f(s)^*$. A three-$\rho$ term plus its complex conjugate is therefore odd under $\v k\to-\v k$ in its totally symmetrised part and even in the part produced by the commutators of the ordered charges (Section~\ref{sec:fin_singlet}). For a globally colour-invariant projectile with the symmetrised three-point function~\eqref{eq:A2d} the odd part is the sum of the symmetrised three-$\rho$ terms of Section~\ref{sec:fin_rows}: it is proportional to the three-point function $\sigma$ of the projectile and contains the target only through the C-odd traces with $D^h$, Eq.~\eqref{eq:traceid}. It is given row by row by Eqs.~\eqref{eq:S-I.II-3}, \eqref{eq:S-I.III-3}, \eqref{eq:S-I.V-3}, \eqref{eq:S-II.I-3}, \eqref{eq:S-III.II-3}, \eqref{eq:S-III.III-3}, \eqref{eq:S-III.V-3} and~\eqref{eq:S-III.VI-3}, for a general three-point function $\sigma$ (no model is assumed), and it is odd configuration by configuration of the target; for a weak target it starts with three gluon exchanges. (Special cases, not part of the result: the odd part is absent for a projectile with $\sigma=0$, such as a Gaussian, McLerran--Venugopalan-type projectile, and it vanishes after an average over a C-even target ensemble.) This agrees with the classical-field result that the leading-order spectrum is even configuration by configuration and that the odd part comes from the term cubic in the projectile charge, averages out for a Gaussian projectile and is ``defined by odderon exchanges''~\cite{McLerran:2016snu}, and with the counting of~\cite{Kovchegov:2018jun}, three scatterings in the projectile and three in the target. The even commutator part does \emph{not} vanish for any projectile: the three-$\rho$ rows contribute to the averaged spectrum through two-point contact terms $\propto\delta^{(2)}(\v x_i-\v x_j)\,n$, which have no counterpart for classical charges. All colour factors of the averaged terms, the symmetrised parts and the parity statements were verified at the operator level, with charges represented by generators on two or three sites and globally colour-invariant density matrices, including a C-odd state with $\sigma\neq0$ (Section~\ref{sec:fin_singlet}); the operator decomposition of each cubic row into its symmetrised and commutator parts, Eqs.~\eqref{eq:symord}--\eqref{eq:symr21}, was verified with explicit $SU(3)$ charges on three sites for all assignments of the charge points to the sites.

\paragraph{Comparison with Li and Skokov.} The symmetrised three-$\rho$ terms are the counterpart, for non-commuting charges, of the first saturation correction $M_{(1)}^*M_{(3)}+{\rm c.c.}\propto\rho_P^3$ of Li and Skokov~\cite{Li:2021ntm}, whose charges are classical; the contact terms have no classical counterpart. The colour structures agree: their two-Wilson-line term has a vanishing $d^{abc}$ projection, as the corresponding structures here, so that this sector does not contribute under (A2); the odd part comes from three Wilson lines and is C-odd in the target; and in our result the invariants with four Wilson lines cancel between Rows~II.I and~III.II in the sum of the symmetrised terms (their kernels are equal and opposite, $\mathfrak K^{(1)}_{\rm II.I}=-\mathfrak K_{\rm III.II}$), so that the odd part has at most three Wilson lines, as in~\cite{McLerran:2016snu,Li:2021ntm}. The parities in $\v k$ do not agree, and this is an open point. Li and Skokov find also a part of $M_{(1)}^*M_{(3)}$ that is even in $\v k$ configuration by configuration (their two-line term and the real part of their three-line kernel), whose $d^{abc}$ projection is not removed by any symmetry, whereas our symmetrised cubic terms are purely odd. Conversely, our result does not reproduce only their odd part: in the two-Wilson-line sector, where no vanishing energy denominator occurs, their term (from the classical non-linear Weizs\"acker--Williams field) is even in $\v k$, while ours (the cubic part of Row~I.II) is odd and has a different kernel. The difference can be traced to a relative phase $i$. The transverse kernel of the classical $O(g^3)$ Weizs\"acker--Williams field of~\cite{Li:2021ntm} occurs in our wave function, as the coefficient of $\log(\Lambda/k^+)$ in the two-$\rho$ one-gluon component, the first term of Eq.~\eqref{eq:A3tworho} (to $10^{-12}$), but with a factor $i$ relative to the LO coefficient~\eqref{eq:PsiLO}, whereas classically the $O(g^3)$ field is in phase with the $O(g)$ one. In the same way the two-charge structures $[\mf A^{(3)}]_{\rho\rho}$ (Row~I.II), the $f\rho$ part of $\mf B_2$ relative to its $\{\rho,\rho\}$ part (Rows~II.I, III.II) and the rescattering of a produced gluon on the rotated second charge after the collision ($\mf B_3$, Row~I.III) carry a phase that differs by $i$ from the classical field and from LCPT with retarded propagators~\cite{Kovchegov:2018jun}, while for the merging of two produced gluons after the collision ($\mf C_1$) the assignments agree. A possible origin is the prescription for $1/k^+$, i.e.\ the $x^-$ profile of the field (retarded, as in the classical solution, or symmetric, which corresponds to real LO coefficients in momentum space); we have not checked this. The $\v k$-even part of~\cite{Li:2021ntm} has not been confirmed by an independent calculation either, and a kernel-level comparison of the odd three-line parts, which requires all cubic rows in momentum space, has not been carried out. The two-line difference is invisible in the averaged spectrum under (A2), since both two-line terms vanish after the $d^{abc}$ projection, but it is present configuration by configuration and in correlations.

\paragraph{Dilute target.} For $U=1+O(\alpha)$ the LO spectrum reduces to the $k_T$-factorised form, Eq.~\eqref{eq:kTfact} below, with the Lipatov-type vertex. In this limit the BFKL and JIMWLK terms reduce to the LO BFKL evolution of the projectile and of the target unintegrated gluon distributions, respectively. The comparison of the NLO terms with $k_T$ factorisation is the subject of Section~\ref{sec:check_dilute}.

\subsection{The dilute--dilute limit}\label{sec:check_dilute}

\paragraph{The limit.} In momentum space the LO amplitude~\eqref{eq:ALOexpl} is
\begin{align}\label{eq:ALOmom}
\mathcal A^{a}_{{\rm LO}\,i}(k)=\frac{\sqrt2\,g}{\sqrt{k^+}}\int\frac{d^2q}{(2\pi)^2}\,\Big[\frac{k^i}{\v k^2}-\frac{q^i}{\v q^2}\Big]\,U^{ab}(\v q-\v k)\,\rho^b(-\v q)\,,\qquad
U^{ab}(\v l)=\int_xe^{i\v l\cdot\v x}U^{ab}(\v x)\,,
\end{align}
so that the gluon is produced with the Lipatov-type vertex $\v k/\v k^2-\v q/\v q^2$, where $\v q$ is the momentum supplied by the projectile and $\v k-\v q$ the momentum transferred by the target. We write $U^{ab}=\delta^{ab}+u^{ab}$, $u^{ab}=f^{acb}\alpha^c$, keep $O(u)$ in the amplitude and in the conjugate amplitude (two-gluon exchange with the target) and the two-point function of the projectile (one gluon exchanged with the projectile on each side); since the charges do not commute, products of two or three charges contribute to this order through $[\rho^a(\v x),\rho^b(\v y)]=if^{abc}\rho^c(\v x)\delta^{(2)}(\v x-\v y)$, while the symmetrised three-point function $\sigma$ is dropped. With $\langle\alpha^c(\v l)\alpha^d(-\v l')\rangle=\delta^{cd}(2\pi)^2\delta^{(2)}(\v l-\v l')T(\v l)$ and $\langle\rho^b(\v q')\rho^{b'}(-\v q)\rangle=\delta^{bb'}(2\pi)^2\delta^{(2)}(\v q-\v q')P(\v q)$ and
\begin{align}\label{eq:dilute_not}
\mathcal N_{\rm d}\equiv\frac{\alpha_sN_c(N_c^2-1)}{\pi^2}\,\frac{S_\perp}{k^+}\,,\qquad\phi_P(\v q)\equiv\frac{P(\v q)}{\v q^2}\,,\qquad\phi_T(\v l)\equiv\v l^2\,T(\v l)\,,
\end{align}
the unintegrated gluon distributions of projectile and target, and $S_\perp$ the transverse area, the LO spectrum takes the $k_T$-factorised form~\cite{Gribov:1983ivg,Kovchegov:2001sc,Kharzeev:2003wz,Blaizot:2004wu}
\begin{align}\label{eq:kTfact}
\frac{d^3N_{\rm LO}}{d^3k}=\mathcal N_{\rm d}\int\frac{d^2q}{(2\pi)^2}\,\frac{\phi_P(\v q)\,\phi_T(\v k-\v q)}{\v k^2}\,.
\end{align}

\paragraph{Real emission (Groups II and III).} At $O(u)$ the anticommutator parts of the two-gluon amplitude~\eqref{eq:A2gamp} cancel, $\mathcal A_B^{(2)}+\mathcal A^{(1)}_{AA}+\mathcal A^{(2)}_{AA}=\frac14[L_k,\mf A_p]+\frac14[L_p,\mf A_k]+O(u^2)$, with $L$ the LO amplitude before the Fourier transform, so that the two-gluon amplitude is linear in $\rho$, and with $\v q_a$ the momentum from the projectile, $\v q_b=\v k+\v p-\v q_a$ that from the target, $x=k^+/(k^++p^+)$,
\begin{align}\label{eq:dilute_real}
\frac{d^3N^{\rm (II+III)}}{d^3k}\bigg|_{\rm dilute}=\mathcal N_{\rm d}\int\frac{d^2q_a}{(2\pi)^2}\int_\Lambda^\vee\frac{dp^+}{p^+}\int d^2p\,\frac{\alpha_sN_c}{\pi^2}\,\phi_P(\v q_a)\,\phi_T(\v q_b)\,\frac{A_{gg}(\v q_a,\v q_b;\v k,\v p;x)}{\v q_a^2\,\v q_b^2}\,,
\end{align}
where $A_{gg}$ is the Born cross section for the production of two gluons in quasi-multi-Regge kinematics~\cite{Fadin:1998py}, in the form given in Ref.~\cite{Bartels:2006hg} (their $A_{\rm gluons}$). We verified this identity numerically at random kinematics for all $0<x<1$, for $SU(2)$ and $SU(3)$, and for $d_\perp$ transverse polarisations ($d_\perp=1.7,3,6$) against the $(D-2)$ terms of~\cite{Bartels:2006hg} (relative deviation $2\times10^{-15}$; there is no free normalisation). It requires the commutator (contact) parts generated by the independent-emission terms, $\mf B_2$ with $\{\rho,\rho\}$ and $\mathcal A^{(1,2)}_{AA}$: without them, or without $\mathcal A_{CA}$, the ratio deviates from one by factors $0.2$--$10$. In the limits $p^+\gg k^+$ and $p^+\ll k^+$, $A_{gg}\to\v q_a^2\v q_b^2/(\v k^2\v p^2)$: these are the real parts of one step of LO BFKL evolution of $\phi_P$ over $\log(\vee/k^+)$ and of $\phi_T$ over $\log(k^+/\Lambda)$, with $\bar\alpha_s=\alpha_sN_c/\pi$; the virtual parts are those of Sections~\ref{sec:check_BFKL} and~\ref{sec:check_JIMWLK}. The coefficient of $\log\mu^2$ and the scale of the coupling, $1-\frac{\alpha_s\beta_0}{4\pi}\log\frac{\v k^2}{\mu^2}$ in the local form~\eqref{eq:dg_local}, agree with~\cite{Fadin:1998py,Bartels:2006hg}. Since $dp^+/p^+=dx/(x(1-x))$ is, at fixed $k^+$, minus the rapidity measure of the second gluon, the finite part of Eq.~\eqref{eq:dilute_real} (the $x$ integral of $A_{gg}/(\v q_a^2\v q_b^2)-1/(\v k^2\v p^2)$) coincides identically with the two-gluon contribution minus its multi-Regge limit of~\cite{Bartels:2006hg} for $s_\Lambda\to\infty$ (numerically at three points, to $10^{-6}$ at $s_\Lambda=10^6$, the difference falling like $1/s_\Lambda$): for the real emission the plus-momentum cut-offs require no conversion, at fixed $\v p\neq0$. For $|\v p|\lesssim\Lambda|\v k|/k^+$ the cut $p^+>\Lambda$ removes part of the soft region at the rapidity of the measured gluon, and the remainder is not uniform there. The two rapidity logarithms refer to plus-momentum cut-offs; on the projectile side this is the asymmetric choice $s_0=\v k_J^2$ of~\cite{Bartels:2006hg}, on the target side they use the minus-momentum fraction. The two schemes differ by $\log(\v p^2/\v k^2)$-weighted LO kernels and trajectories (energy-scale terms, not constants times LO); for the real emission the shift cancels between the two ends, for the virtual trajectories (split of the loop-gluon window at $p^+=k^+$ instead of $y_p=y_k$) it has not been evaluated.

% DILUTE-SLOT BEGIN --------------------------------------------------------------------------
% Status of the comparison of the virtual correction with the one-loop central-emission vertex, as verified in
% notes/ref3_dilute.tex. The result of the completed comparison is to replace the text between the two markers.
\paragraph{Virtual correction (Group I) for a point-like quark.} For $\rho^a(\v x)=t^a\delta^{(2)}(\v x)$, Row~I.I and the two- and three-$\rho$ parts of $\mf A^{(3)}$ vanish ($f^{abc}\{t^b,t^c\}=0$, scaleless integrals), and we evaluated the remaining Rows~I.II (one-$\rho$ part), I.III (principal value, cut-offs), I.IV and I.V in momentum space, in $d_\perp=2-\epsilon$, $\MS$, at $\mu=|\v k|$, at two kinematic points. To the accuracy of the numerical integration ($2\times10^{-4}$ in the coefficients of the logarithms, $2\times10^{-3}$ in the constant),
\begin{align}\label{eq:dilute_groupI}
\frac{\text{Group I}}{\rm LO}\bigg|_{\rm point\ quark}&=2\omega(\v q_a)\log\frac{\vee}{k^+}+2\omega(\v q_b)\log\frac{k^+}{\Lambda}\notag\\&\quad+\frac{\alpha_sN_c}{\pi}\Big[\frac2\epsilon\log\frac{k^+}{\Lambda}+\log^2\frac{k^+}{\Lambda}\Big]+C(\v k,\v l)\,,\notag\\
2\omega(\v q)&\equiv\frac{\alpha_sN_c}{\pi}\Big[\frac2\epsilon-\log\frac{\v q^2}{\mu^2}\Big],
\end{align}
with $\v q_a=\v k+\v l$, $\v q_b=-\v l$ ($-\v l$ the momentum from the target); $\omega$ has the form of the one-loop gluon trajectory of~\cite{Fadin:1998py} in our normalisation. (i) The coefficient of $\log(\vee/k^+)$ is $2\omega(\v q_a)$ exactly, pole and finite part; it comes from Rows~I.III and I.IV. (ii) The coefficient of $\log(k^+/\Lambda)$ is $2\omega(\v q_b)$, at a kinematic point where $2\omega(\v q_a)$ is very different, but only in the sum of Rows~I.II, I.III, I.IV and I.V (Row~I.II alone has $\omega(\v k)$ and $\omega(\v q_a)$ on the two terms of the vertex); on top of it there is the $\v l$-independent soft--collinear term $\frac{\alpha_sN_c}{\pi}[\frac2\epsilon L+L^2]$, $L=\log(k^+/\Lambda)$, whose dependence on $\mu$ and $\v k^2$ was not separated. (iii) The $\frac{11}{3}$ poles cancel between Rows~I.II and~I.V, and the unresolved part of Row~I.V cancels the whole self-energy part of Row~I.II, $\xi$ by $\xi$ to $10^{-13}$, including the $\log(\xi(1-\xi))$ that produces $\frac{67}{9}$: Group~I contains no $\frac{67}{9}$ (Eq.~\eqref{eq:dg_SEid}). (iv) The kinematics-dependent double logarithms cancel between Row~I.III and the resolved part of Row~I.V; the remaining $\frac{\alpha_sN_c}{\pi}L^2$ is part of the soft--collinear term, so that Group~I alone is not free of $\log^2(k^+/\Lambda)$. (v) The finite remainder $C(\v k,\v l)$ is a function, not a constant: $0.055$ and $0.305$ times LO at the two points. (vi) In this limit the result does not depend on the prescription for the vanishing energy denominators: the on-shell part of $\mf B_3$ is purely imaginary relative to LO and drops out.

\paragraph{What is and what is not tested.} The one-loop central-emission vertex~\cite{Fadin:1998py,DelDuca:1998cx,Bartels:2006hg} regulates the soft--collinear region by $\epsilon^{-2}(\boldsymbol\Delta^2/\mu^2)^\epsilon$, $\boldsymbol\Delta=\v q_a+\v q_b=\v k$, and has no rapidity cut-off; in Eq.~\eqref{eq:dilute_groupI} the same region is regulated by the plus-momentum cut-off. The two virtual parts therefore differ already in their pole terms by a function of $\boldsymbol\Delta$ that is not of the form $F_P(\v q_a)+F_T(\v q_b)$, the form that at fixed order could be absorbed into $\phi_P$ and $\phi_T$; the scheme-independent comparison requires the virtual and the real parts with matched regulators, including the real emission at $|\v p|\lesssim\Lambda|\v k|/k^+$ with the cut $p^+>\Lambda$, which in dimensional regularisation produces $\Lambda$-dependent soft--collinear terms of the type $\frac2\epsilon L$ and $L^2$. This comparison has not been carried out: whether the extra term of Eq.~\eqref{eq:dilute_groupI} cancels against that real contribution, and whether the finite remainder agrees with the one-loop vertex, is open. No disagreement with~\cite{Fadin:1998py,Bartels:2006hg} was found. The tests of this section are made at the level of the amplitudes and coefficients, Eqs.~\eqref{eq:A2gamp}, \eqref{eq:B2mom}, \eqref{eq:C1X} and~\eqref{eq:A1X}; the dilute limit of the coordinate-space finite terms as printed, and the equality of the printed colour-singlet contact terms with the commutator contributions required by Eq.~\eqref{eq:dilute_real}, have not been checked. For a general (not point-like) projectile Group~I has not been reduced in this limit.
% DILUTE-SLOT END ----------------------------------------------------------------------------

\subsection{Projectile evolution: the BFKL term}\label{sec:check_BFKL}

\paragraph{Internal consistency.} The ordered form~\eqref{eq:BFKLrho} is an exact regrouping of the term-by-term list of all two-$\rho$ terms proportional to $\log(\vee/k^+)$, for an arbitrary ordered correlator (relative difference $\lesssim10^{-15}$ for $SU(2)$, $10^{-14}$ for $SU(3)$), and Eq.~\eqref{eq:BFKLsinglet} is its colour-singlet reduction to the same accuracy.

\paragraph{Exact one-gluon evolution of a colour-singlet projectile.} As an independent test we computed the change of the spectrum due to one additional gluon with $k^+<p^+<\vee$ in a colour-singlet projectile directly, without the expansion in $\rho$ used in the rest of the paper. Any set of colour sources $s$ at positions $\v x_s$, with charges $Q_s^e$, produces the measured gluon with the operator $\sum_{s,s'}Q^e_{s'}Q^e_sF(\v x_s,\v x_{s'})$, where
\begin{align}
F(\v a,\v b)=\int_{w,w'}e^{-i\v k\cdot(\v w-\v w')}\Kv(\v a-\v w)\cdot\Kv(\v b-\v w')\,{\rm Tr}\big[(U^\dagger(\v a)-U^\dagger(\v w))(U(\v b)-U(\v w'))\big]\,,
\end{align}
since the colour rotation of the sources cancels in the inclusive sum. For a quark--antiquark dipole at $\v x_1$, $\v x_2$ the three charges of the $q\bar qg$ colour singlet satisfy $Q_1^2=Q_2^2=C_F$, $Q_g^2=N_c$, $Q_1\cdot Q_2=1/(2N_c)$ and $Q_{1,2}\cdot Q_g=-N_c/2$, and the probability to emit the additional gluon at $\v z$ is $\frac{g^2C_F}{4\pi^3}\log\frac{\vee}{k^+}K_D(\v x_1,\v x_2,\v z)=\frac{\alpha_sC_F}{\pi^2}\log\frac{\vee}{k^+}K_D(\v x_1,\v x_2,\v z)$, i.e.\ $-2\mf N_2$ of Eq.~\eqref{eq:N2X} for the dipole. With the LO spectrum of the dipole, $\frac{d^3N_{\rm LO}}{d^3k}=\frac{\alpha_s}{4\pi^4k^+}\mathcal D_{12}$ by Eq.~\eqref{eq:LOsinglet}, real minus virtual gives, exactly at finite $N_c$,
\begin{align}
\delta\Big[\frac{d^3N}{d^3k}\Big]&=\frac{\alpha_sN_c}{2\pi^2}\log\frac{\vee}{k^+}\int_zK_D(\v x_1,\v x_2,\v z)\big[\mathcal D_{1z}+\mathcal D_{z2}-\mathcal D_{12}\big]\times\frac{\alpha_s}{4\pi^4k^+}\,,\notag\\
\mathcal D_{ab}&=\frac{1}{2N_c}\big[F_{aa}+F_{bb}-F_{ab}-F_{ba}\big]\,,
\end{align}
the dipole BFKL (linearised BK) kernel acting on the projectile. Substituting $n=C_F[\delta_{11}+\delta_{22}-\delta_{12}-\delta_{21}]$ into Eq.~\eqref{eq:BFKLsinglet} reproduces this result pointwise in $\v z$. Since every symmetric $n(\v x,\v y)$ with vanishing row sums, i.e.\ every colour-singlet projectile, is a linear combination of such dipole correlators, and Eq.~\eqref{eq:BFKLsinglet} is linear in $n$, the BFKL term is the exact leading-logarithmic projectile evolution for every colour-singlet projectile. The calculation with fundamental charges is exact to all orders in the charges, which shows that the expansion in $\rho$ loses nothing for a dipole.

\paragraph{Literature.} This is the structure found by Kovner and Lublinsky for the single inclusive spectrum~\cite{Kovner:2006wr} and by Kovner, Lublinsky and Weigert~\cite{Kovner:2006ge}: the projectile correlator evolves with BFKL from the projectile rapidity down to the rapidity of the produced gluon. It is the first-order expansion of the leading-logarithmic factorisation theorem of Gelis, Lappi and Venugopalan~\cite{Gelis:2008rw}, and it satisfies the general structure of the NLO correction discussed by Gelis and Laidet~\cite{Gelis:2012ri}.

\subsection{Cubic terms on exact charge states: a quark and a dipole}\label{sec:check_cubic}

For a single quark, $\rho^a(\v x)=t^a\delta^{(2)}(\v x-\v x_0)$ in a colour-averaged state, the three-point function is exact: $\langle\rho^a\rho^b\rho^c\rangle=\frac1{N_c}{\rm tr}(t^at^bt^c)=\frac{1}{4N_c}(d^{abc}+if^{abc})$ at $\v x_0$, so that (A2) holds with $\sigma=\frac{(N_c^2-4)(N_c^2-1)}{4N_c^2}=C_F\frac{N_c^2-4}{2N_c}$ ($\frac{10}{9}$ for $SU(3)$; $d^{abc}t^at^bt^c=C_F\frac{N_c^2-4}{2N_c}$) and $n=C_F$. Each cubic colour structure then has the unique value $\frac{1}{4N_c}[\Delta+i\Phi]$ for an ordered product and $\frac{1}{4N_c}\Delta$ for a product with an anticommutator (${\rm tr}\,t^a\{t^b,t^c\}=\frac12d^{abc}$), with $\Delta=\sum C^{abc}d^{abc}$ and $\Phi=C^{abc}f^{abc}$ at $\v x_0$. Eq.~\eqref{eq:cubicavg} reproduces this term by term: the symmetrised part is the $d^{abc}$ half, the contact part the $f^{abc}$ half, and for an antiquark the first changes sign and the second does not. For a colour-singlet dipole, quark at $\v x_0$ and antiquark at $\v x_1$, $n=C_F\,s_is_j$ and $\sigma=\sigma_q\,s_is_js_k$, with $s=+1$ at the quark and $-1$ at the antiquark; the symmetrised part of a cubic term carries $s_is_js_k$ and is the part odd under the exchange of quark and antiquark, i.e.\ under the C-conjugation of the dipole, while the contact part carries products of two $s$'s and is even. The $\v k$-odd cubic term of a dipole is therefore the part of the spectrum that distinguishes the quark end from the antiquark end, $\frac12[\sigma(\v x_0,\v x_1)-\sigma(\v x_1,\v x_0)]$; it needs a C-odd target and vanishes for an orientation-averaged dipole, whereas the $\v k$-even cubic (contact) part does not. We evaluated all thirteen cubic colour structures with explicit $SU(2)$, $SU(3)$ and $SU(4)$ matrices for the quark, the antiquark, the dipole and the conjugate dipole, all assignments of the charge points, and random Wilson lines: the exact ordered expectation values agree with the split~\eqref{eq:cubicavg} to $5\times10^{-14}$, and the C-split $\frac12[\sigma(P)\mp\sigma(P^{\rm C})]$ to $8\times10^{-15}$. For a quark or a dipole neither part is enhanced: both are of the order of the quadratic terms (colour weights relative to the LO weight for a dense random $SU(3)$ target: on average $0.12$ (range $0.05$--$0.23$ over the structures) for the symmetrised and $0.8$ (range $0.4$--$1.1$) for the contact part; for a weak target the symmetrised part is suppressed by one power of the target field). For a projectile of $m$ independent dipoles $n$ and $\sigma$ are both proportional to $m$. The language of a ``first saturation correction'' therefore applies to the symmetrised part only for a coherent dense projectile with large $\sigma$; this is a property of the projectile model, not of the calculation. The kernels at the triple coincidence of the charge points, which a single quark requires, are integrable (Appendix~\ref{sec:fin_integrable}).

\subsection{Target evolution: the JIMWLK term}\label{sec:check_JIMWLK}

We computed $-H_{\rm JIMWLK}\,\mathcal O_{\rm LO}$ independently, with the Hamiltonian~\eqref{eq:HJIMWLK} in the normalisation of~\cite{Lublinsky:2016meo,Kovner:2013ona}, for a general ordered correlator. The term-by-term list of the two-$\rho$ terms proportional to $\log(k^+/\Lambda)$, the grouped form~\eqref{eq:JIMWLKrho} and $-H_{\rm JIMWLK}$ acting on the LO spectrum agree to $10^{-14}$ for $SU(2)$ and $SU(3)$, and the colour-singlet form~\eqref{eq:JIMWLKsinglet} agrees with the colour-singlet projection to $10^{-15}$. The JIMWLK term is therefore the JIMWLK evolution of the LO spectrum, with the correct sign and normalisation, for an arbitrary projectile and not only for a colour-singlet one. Together with Section~\ref{sec:check_BFKL}, this confirms at fixed order the factorisation of the rapidity logarithms at the rapidity of the produced gluon~\cite{Kovner:2006wr,Kovner:2006ge,Gelis:2008rw}.

\subsection{Wave function}\label{sec:check_wf}

The components of the dressed vacuum, which are quoted in Section~\ref{sec:vacin} from Lublinsky and Mulian~\cite{Lublinsky:2016meo}, correspond to their equations as follows:
\begin{center}
\renewcommand{\arraystretch}{1.2}
\small
\begin{tabular}{@{}p{4.6cm}p{2.3cm}p{6.6cm}@{}}
\toprule
component & Ref.~\cite{Lublinsky:2016meo} & comparison\\
\midrule
LO soft gluon, normalisation $\mf N_2$ & eqs.~2.50--2.52 & identical\\
two gluons, $O(g^2)$ & eqs.~3.95--3.96 & identical\\
one gluon, $O(g^3)$, linear in $\rho$ & eqs.~3.97--3.98 & identical for $N_f=0$, incl.\ $\frac{11}{3}$, $\frac{67}{9}-\frac{\pi^2}{3}$, $\log^2$ cut-off terms\\
one gluon, $O(g^3)$, $f^{abc}\{\rho^b,\rho^c\}$ & eqs.~3.100--3.101 & equal after the antisymmetrisation enforced by $f^{abc}\{\rho^b,\rho^c\}$\\
one gluon, $O(g^3)$, $\rho\rho\rho$ & eq.~3.102 & identical up to the operator ordering, see below\\
\bottomrule
\end{tabular}
\end{center}
In the two-$\rho$ component the two forms differ by terms symmetric under the exchange of the two charges, which do not contribute. The operator ordering of the three-$\rho$ component, $\rho^a(\v w)\rho^b(\v x)\rho^b(\v y)$, i.e.\ $[\mf A^{(3)}]_{\rho\rho\rho}=\mf A\,\mf N_2$, follows from third-order perturbation theory with operator-valued matrix elements; we verified it against a construction of $\Omega$ and an exact diagonalisation with finite-dimensional matrices representing the charges (Appendix~\ref{app:pt_num}; the opposite ordering deviates at order one). The components of the dressed one-gluon state of Section~\ref{sec:onegin} were rebuilt numerically from the matrix elements~\eqref{eq:m1}--\eqref{eq:m5} and Eq.~\eqref{eq:pt2}, and the unitarity relations of Section~\ref{sec:unitarity} were verified symbolically (Appendix~\ref{app:unitarity}).

\subsection{The constant \texorpdfstring{$\frac{67}{9}-\frac{\pi^2}{3}$}{67/9-pi2/3}}\label{sec:lit_constant}

The term of Row I.II proportional to the LO structure is, with its complex conjugate,
\begin{align}\label{eq:constantLO}
\frac{\alpha_sN_c}{2\pi}\Big(\frac{67}{9}-\frac{\pi^2}{3}\Big)\frac{d^3N_{\rm LO}}{d^3k}\,.
\end{align}
It follows directly from the one-$\rho$ part of $\mf A^{(3)}$, Eq.~\eqref{eq:A3onerho}, and the normalisation of the LO amplitude. Before the cancellation of the ultraviolet $\beta_0$ logarithms against Row I.V, Row I.II with its complex conjugate is $\frac{\alpha_sN_c}{2\pi}\big[\frac{11}{3}\big(2\gamma_E+\frac12\log\frac{X^2\mu^2_{\MS}}{4}+\frac12\log\frac{X'^2\mu^2_{\MS}}{4}\big)+\frac{67}{9}-\frac{\pi^2}{3}\big]\times$LO, $X=x-w$, $X'=x'-w'$. This is, term by term at $N_f=0$, the bare one-gluon kernel $K_{JSJ}$ of the NLO JIMWLK Hamiltonian in the form of Kovner, Lublinsky, Skokov and Zhao~\cite{Kovner:2023vsy}, with the same sign, size and scale; since both are built from the wave function of~\cite{Lublinsky:2016meo}, this is a consistency check rather than an independent comparison. The inclusive quantities in which the same constant appears, the NLO BK kernel of Balitsky and Chirilli~\cite{Balitsky:2007feb}, the NLO BFKL kernel of Fadin and Lipatov~\cite{Fadin:1998py} and the NLO jet vertex of Bartels, Sabio Vera and Schwennsen~\cite{Bartels:2006hg}, have the coefficient $\frac{\alpha_sN_c}{4\pi}(\frac{67}{9}-\frac{\pi^2}{3})$, half of Eq.~\eqref{eq:constantLO}. These quantities are inclusive over the collinear gluon pair; as noted in~\cite{Kovner:2023vsy}, half of the constant sits there in the two-gluon terms, and in~\cite{Fadin:1998py} it appears only after the real two-gluon production is integrated over the pair at fixed produced momentum. The amplitude-level one-loop Lipatov vertex~\cite{DelDuca:1998cx} contains no $\frac{67}{9}$.

\paragraph{The constant is a property of the grouping.} Eq.~\eqref{eq:constantLO} is not a constant of the result. By the identity~\eqref{eq:dg_SEid}, the self-energy part of Row~I.II, which carries the whole constant $\frac{67}{9}$ together with the $\beta_0$ logarithm, is cancelled identically by the part of Row~I.V in which the target rotates the gluon pair at its centre of mass. The constant~\eqref{eq:constantLO} is therefore a property of the grouping of Rows~I.II and~I.V; in the grouping used in this paper (colour of the pair at the position of one of its gluons) the compensating terms are in the genuine two-$\rho$ part of Row~I.V, Eq.~\eqref{eq:S-I.V-2o}. This is confirmed in momentum space in the dilute limit with a point-like projectile, Section~\ref{sec:check_dilute}, where Group~I contains no $\frac{67}{9}$ and its finite part is a function of the kinematics, not a multiple of LO. The value of the LO-proportional constant left in the genuine two-$\rho$ part of Row~I.V in our grouping, for a general projectile and target, has not been evaluated. We do not read Eq.~\eqref{eq:constantLO} as the LO-type constant of the single inclusive spectrum, and we make no statement about where ``the other half'' of the inclusive value resides.

\paragraph{Inclusive combinations.} We have not established that an inclusive combination of our collinear terms reproduces the value of the inclusive kernels, and there is no general reason why it should: in the one-hadron-inclusive gluon coefficient function of $H\to gX$ the constant $\frac{67}{9}$ appears only in the $\delta(1-z)$ part~\cite{Almasy:2011eq}, and the energy-weighted integral is the process-dependent total rate. For Row~II.IV the momentum-weighted ($\int d\xi\,\xi$) collinear terms give, exactly, $\frac{\alpha_sN_c}{4\pi}\big(\frac{65}{18}-\frac{4\pi^2}{3}\big)$ times LO in our scheme (zero in the scheme of CXY), plus a term of linearly polarised structure from the tensor term of Eq.~\eqref{eq:S-II.IVrem}, which is not proportional to LO; the corresponding regions of Rows~II.I, III.III and I.V are not of LO form and have not been integrated. In the dilute limit the real rows are exactly the two-gluon production cross section of quasi-multi-Regge kinematics, Eq.~\eqref{eq:dilute_real}; hence for a pair-inclusive quantity built from our amplitudes (virtual correction at fixed produced momentum plus two-gluon production integrated at fixed pair momentum) the real part contributes the value $\frac{\alpha_sN_c}{4\pi}(\frac{67}{9}-\frac{\pi^2}{3})$ of~\cite{Fadin:1998py}, provided the real rows are integrated with $d_\perp$ polarisations, as in our scheme.

\subsection{Collinear and running-coupling terms}\label{sec:lit_dglap}

The collinear term comes from Row II.IV, in which the measured gluon is produced in the splitting $g\to gg$ of a gluon of momentum $k^+/\xi$ that was produced at LO; it is a final-state splitting. Compared with the forward hadron production results of Chirilli, Xiao and Yuan~\cite{Chirilli:2011km,Chirilli:2012jd} and of Altinoluk et al.~\cite{Altinoluk:2014eka}:
\begin{itemize}
\item The splitting function has the normalisation of the LO $P_{gg}$, $2N_c[\frac{\xi}{(1-\xi)_+}+\frac{1-\xi}{\xi}+\xi(1-\xi)]+\frac{\beta_0}{2}\delta(1-\xi)$ at $N_f=0$, and the coefficient $\beta_0/2$ of $\delta(1-\xi)$ is the one required by the momentum sum rule $\int_0^1d\xi\,\xi P_{gg}(\xi)=0$. The prefactor $\frac{\alpha_s}{2\pi}$, the factor $\xi^{-2}$ with the phase $e^{-i\v k\cdot(\v y-\v y')/\xi}$, the sign and the dependence on $\mu$ are those of a final-state $\MS$ fragmentation of the LO spectrum of the parent gluon.
\item The end point $\xi\to1$ (integrated gluon soft) is plus-subtracted into the JIMWLK term, as the rapidity divergence is subtracted into BK evolution in~\cite{Chirilli:2012jd}. The end point $\xi\to0$ (measured gluon soft with respect to its parent) is subtracted into the BFKL term; the combination of the two subtractions is exactly the ``double plus'' prescription of~\cite{Kovner:2023vsy}, for the same reason: both poles are part of the LL evolution. In the forward (hybrid) case $\xi$ is bounded from below and no subtraction at $\xi\to0$ arises.
\item The argument of the collinear logarithm is the transverse separation of the measured gluon in the amplitude and in the conjugate amplitude, $(\v y-\v y')/\xi$, whereas in the $\MS$ form of~\cite{Chirilli:2012jd,Altinoluk:2014eka} it is the separation of the parent, $\v y-\v y'$. The two forms differ by the finite term $\Delta$ of Eq.~\eqref{eq:dg_Delta}, which results from the $d$-dimensional Jacobian of the $\MS$ counterterm (Section~\ref{sec:dg_msbar}, Appendix~\ref{app:scheme}). Our form coincides with the cut-off scheme of~\cite{Altinoluk:2014eka}, with $\mu^2_{\MS}=c_0^2\mu^2_{\rm cut}$.
\item No counterpart arises to the initial-state DGLAP evolution of the projectile parton distribution of the hybrid formalism: the projectile here is described by an unintegrated charge correlator, whose small-$x$ evolution is the BFKL term.
\item The $\beta_0$ terms. The real emission produces no $\delta(1-\xi)$ term, and the virtual $\frac{11}{3}$ poles of the vertex correction (Row I.II) cancel against those of Row I.V, for any grouping of the finite terms (Section~\ref{sec:dg_poles}); which numbers $\frac{11}{6}$ are $\beta_0/(2N_c)$ is identified there by a BLM-type quark loop (the calculation itself is for $N_f=0$). The term $\frac{\beta_0}{2}\delta(1-\xi)$ of $P_{gg}$ is added and subtracted, the subtracted term being the running-coupling term; with it the LO coupling becomes $g^2(\mu)$ at the scale of the collinear subtraction, or $g^2(c_0^2/r^2)$ with $r$ the transverse separation conjugate to $\v k$. In the classification of~\cite{Kovner:2023vsy}, where of the $2\beta_0$ of the NLO JIMWLK kernel one is the ultraviolet running of the coupling and one a DGLAP-type resolution logarithm, our explicit $\beta_0$ terms are, in the grouping used here, of the second, collinear, type. This agrees in structure with Kovchegov and Weigert~\cite{Kovchegov:2007vf} and with Horowitz and Kovchegov~\cite{Horowitz:2010yg}, in whose calculation (quark--quark scattering, BLM resummation of quark bubbles) the observable includes collinear pairs whose invariant mass, i.e.\ the virtuality of the produced-gluon line, is below a fixed cut-off $\Lambda_{\rm coll}$, and the coupling of the emission vertex of the produced gluon is $\alpha_s(\Lambda^2_{\rm coll}e^{-5/3})$ (Eqs.~(3.10)--(3.16) of~\cite{Horowitz:2010yg}); $\Lambda_{\rm coll}$ does not depend on the exchanged momenta and plays the role of the factorisation scale of the fragmentation function~\cite{Kovchegov:2007vf}, while the remaining couplings run with $\v q^2$, $(\v k-\v q)^2$ and a complex scale $Q(\v q,\v k-\v q)$, Eqs.~(3.31), (3.32) of~\cite{Horowitz:2010yg}. The other couplings of~\cite{Horowitz:2010yg}, associated with the $t$-channel exchanges, are not renormalised in this calculation.
\end{itemize}

\subsection{Finite terms}\label{sec:lit_finite}

\paragraph{Transcription.} The equations of Section~\ref{sec:fin_rows} and of Appendix~\ref{app:finite} were generated mechanically from the checked expressions. Their colour factors as printed were evaluated numerically, with random $SU(3)$ and $SU(4)$ Wilson lines, and compared with the colour tensors of the rows before the factorisation into differences of Wilson lines and, for the averaged form, with the direct contractions of these tensors with $d^{abc}$, with $f^{abc}$ at coincident charges and with $\delta^{ab}$ (largest relative deviation $7\times10^{-14}$); the trace identities~\eqref{eq:traceid}, the symmetrisation identities~\eqref{eq:symord}--\eqref{eq:symr21} at the operator level, and the rewriting of the longitudinal integrals in terms of the momentum fraction $\xi$, by quadrature, were checked in the same way.

\paragraph{Soft remainders.} The soft remainders of the plus prescriptions cancel in four pairs of a virtual and a real term (Section~\ref{sec:nuance_soft}, Table~\ref{tab:pairs}). The real-minus-virtual combination is the analogue of the remainder $[A(y-z)-A(\bar y-z)]^2[s(y,z)s(z,\bar y)-s(y,\bar y)]$ of the kinematically constrained forward production of~\cite{Altinoluk:2014eka}, and the soft remainders can equivalently be absorbed into a JIMWLK term in which the soft gluon is restricted to rapidities below that of the measured gluon, Eq.~\eqref{eq:kinconstr}, as in the kinematically constrained formulations of~\cite{Altinoluk:2014eka,Ducloue:2016shw} (the constraint of~\cite{Ducloue:2016shw} has not been compared in detail). The analogous remainders of forward production were discussed in~\cite{Stasto:2013cha,Altinoluk:2014eka,Ducloue:2016shw}.

